\documentclass[
  doc,
  floatsintext,
  longtable,
  a4paper,
  nolmodern,
  notxfonts,
  notimes,
  12pt,
  colorlinks=true,linkcolor=blue,citecolor=blue,urlcolor=blue]{apa7}

\usepackage{amsmath}
\usepackage{amssymb}

\geometry{inner=1in, outer=1in}
\fancyhfoffset[LE,RO]{0cm}


\usepackage[bidi=default]{babel}
\babelprovide[main,import]{spanish}


% get rid of language-specific shorthands (see #6817):
\let\LanguageShortHands\languageshorthands
\def\languageshorthands#1{}

\RequirePackage{longtable}
\RequirePackage{threeparttablex}

\makeatletter
\renewcommand{\paragraph}{\@startsection{paragraph}{4}{\parindent}%
	{0\baselineskip \@plus 0.2ex \@minus 0.2ex}%
	{-.5em}%
	{\normalfont\normalsize\bfseries\typesectitle}}

\renewcommand{\subparagraph}[1]{\@startsection{subparagraph}{5}{0.5em}%
	{0\baselineskip \@plus 0.2ex \@minus 0.2ex}%
	{-\z@\relax}%
	{\normalfont\normalsize\bfseries\itshape\hspace{\parindent}{#1}\textit{\addperi}}{\relax}}
\makeatother




\usepackage{longtable, booktabs, multirow, multicol, colortbl, hhline, caption, array, float, xpatch}
\usepackage{subcaption}


\renewcommand\thesubfigure{\Alph{subfigure}}
\setcounter{topnumber}{2}
\setcounter{bottomnumber}{2}
\setcounter{totalnumber}{4}
\renewcommand{\topfraction}{0.85}
\renewcommand{\bottomfraction}{0.85}
\renewcommand{\textfraction}{0.15}
\renewcommand{\floatpagefraction}{0.7}

\usepackage{tcolorbox}
\tcbuselibrary{listings,theorems, breakable, skins}
\usepackage{fontawesome5}

\definecolor{quarto-callout-color}{HTML}{909090}
\definecolor{quarto-callout-note-color}{HTML}{0758E5}
\definecolor{quarto-callout-important-color}{HTML}{CC1914}
\definecolor{quarto-callout-warning-color}{HTML}{EB9113}
\definecolor{quarto-callout-tip-color}{HTML}{00A047}
\definecolor{quarto-callout-caution-color}{HTML}{FC5300}
\definecolor{quarto-callout-color-frame}{HTML}{ACACAC}
\definecolor{quarto-callout-note-color-frame}{HTML}{4582EC}
\definecolor{quarto-callout-important-color-frame}{HTML}{D9534F}
\definecolor{quarto-callout-warning-color-frame}{HTML}{F0AD4E}
\definecolor{quarto-callout-tip-color-frame}{HTML}{02B875}
\definecolor{quarto-callout-caution-color-frame}{HTML}{FD7E14}

%\newlength\Oldarrayrulewidth
%\newlength\Oldtabcolsep


\usepackage{hyperref}



\usepackage{color}
\usepackage{fancyvrb}
\newcommand{\VerbBar}{|}
\newcommand{\VERB}{\Verb[commandchars=\\\{\}]}
\DefineVerbatimEnvironment{Highlighting}{Verbatim}{commandchars=\\\{\}}
% Add ',fontsize=\small' for more characters per line
\usepackage{framed}
\definecolor{shadecolor}{RGB}{241,243,245}
\newenvironment{Shaded}{\begin{snugshade}}{\end{snugshade}}
\newcommand{\AlertTok}[1]{\textcolor[rgb]{0.68,0.00,0.00}{#1}}
\newcommand{\AnnotationTok}[1]{\textcolor[rgb]{0.37,0.37,0.37}{#1}}
\newcommand{\AttributeTok}[1]{\textcolor[rgb]{0.40,0.45,0.13}{#1}}
\newcommand{\BaseNTok}[1]{\textcolor[rgb]{0.68,0.00,0.00}{#1}}
\newcommand{\BuiltInTok}[1]{\textcolor[rgb]{0.00,0.23,0.31}{#1}}
\newcommand{\CharTok}[1]{\textcolor[rgb]{0.13,0.47,0.30}{#1}}
\newcommand{\CommentTok}[1]{\textcolor[rgb]{0.37,0.37,0.37}{#1}}
\newcommand{\CommentVarTok}[1]{\textcolor[rgb]{0.37,0.37,0.37}{\textit{#1}}}
\newcommand{\ConstantTok}[1]{\textcolor[rgb]{0.56,0.35,0.01}{#1}}
\newcommand{\ControlFlowTok}[1]{\textcolor[rgb]{0.00,0.23,0.31}{\textbf{#1}}}
\newcommand{\DataTypeTok}[1]{\textcolor[rgb]{0.68,0.00,0.00}{#1}}
\newcommand{\DecValTok}[1]{\textcolor[rgb]{0.68,0.00,0.00}{#1}}
\newcommand{\DocumentationTok}[1]{\textcolor[rgb]{0.37,0.37,0.37}{\textit{#1}}}
\newcommand{\ErrorTok}[1]{\textcolor[rgb]{0.68,0.00,0.00}{#1}}
\newcommand{\ExtensionTok}[1]{\textcolor[rgb]{0.00,0.23,0.31}{#1}}
\newcommand{\FloatTok}[1]{\textcolor[rgb]{0.68,0.00,0.00}{#1}}
\newcommand{\FunctionTok}[1]{\textcolor[rgb]{0.28,0.35,0.67}{#1}}
\newcommand{\ImportTok}[1]{\textcolor[rgb]{0.00,0.46,0.62}{#1}}
\newcommand{\InformationTok}[1]{\textcolor[rgb]{0.37,0.37,0.37}{#1}}
\newcommand{\KeywordTok}[1]{\textcolor[rgb]{0.00,0.23,0.31}{\textbf{#1}}}
\newcommand{\NormalTok}[1]{\textcolor[rgb]{0.00,0.23,0.31}{#1}}
\newcommand{\OperatorTok}[1]{\textcolor[rgb]{0.37,0.37,0.37}{#1}}
\newcommand{\OtherTok}[1]{\textcolor[rgb]{0.00,0.23,0.31}{#1}}
\newcommand{\PreprocessorTok}[1]{\textcolor[rgb]{0.68,0.00,0.00}{#1}}
\newcommand{\RegionMarkerTok}[1]{\textcolor[rgb]{0.00,0.23,0.31}{#1}}
\newcommand{\SpecialCharTok}[1]{\textcolor[rgb]{0.37,0.37,0.37}{#1}}
\newcommand{\SpecialStringTok}[1]{\textcolor[rgb]{0.13,0.47,0.30}{#1}}
\newcommand{\StringTok}[1]{\textcolor[rgb]{0.13,0.47,0.30}{#1}}
\newcommand{\VariableTok}[1]{\textcolor[rgb]{0.07,0.07,0.07}{#1}}
\newcommand{\VerbatimStringTok}[1]{\textcolor[rgb]{0.13,0.47,0.30}{#1}}
\newcommand{\WarningTok}[1]{\textcolor[rgb]{0.37,0.37,0.37}{\textit{#1}}}

\providecommand{\tightlist}{%
  \setlength{\itemsep}{0pt}\setlength{\parskip}{0pt}}
\usepackage{longtable,booktabs,array}
\newcounter{none} % for unnumbered tables
\usepackage{calc} % for calculating minipage widths
% Correct order of tables after \paragraph or \subparagraph
\usepackage{etoolbox}
\makeatletter
\patchcmd\longtable{\par}{\if@noskipsec\mbox{}\fi\par}{}{}
\makeatother
% Allow footnotes in longtable head/foot
\IfFileExists{footnotehyper.sty}{\usepackage{footnotehyper}}{\usepackage{footnote}}
\makesavenoteenv{longtable}

\usepackage{graphicx}
\makeatletter
\newsavebox\pandoc@box
\newcommand*\pandocbounded[1]{% scales image to fit in text height/width
  \sbox\pandoc@box{#1}%
  \Gscale@div\@tempa{\textheight}{\dimexpr\ht\pandoc@box+\dp\pandoc@box\relax}%
  \Gscale@div\@tempb{\linewidth}{\wd\pandoc@box}%
  \ifdim\@tempb\p@<\@tempa\p@\let\@tempa\@tempb\fi% select the smaller of both
  \ifdim\@tempa\p@<\p@\scalebox{\@tempa}{\usebox\pandoc@box}%
  \else\usebox{\pandoc@box}%
  \fi%
}
% Set default figure placement to htbp
\def\fps@figure{htbp}
\makeatother







\usepackage{newtx}

\defaultfontfeatures{Scale=MatchLowercase}
\defaultfontfeatures[\rmfamily]{Ligatures=TeX,Scale=1}





\title{Nombres de variables}


\shorttitle{NOMBRES DE VARIABLES}


\usepackage{etoolbox}



\ccoppy{\textcopyright~2026}



\author{Edison Achalma}



\affiliation{
{Escuela Profesional de Economía, Universidad Nacional de San Cristóbal
de Huamanga}}




\leftheader{Achalma}

\date{2026-10-20}


\abstract{Este abstract será actualizado una vez que se complete el
contenido final del artículo. }



\authornote{\par{\addORCIDlink{Edison Achalma}{0000-0001-6996-3364}} 
\par{ }
\par{   El autor no tiene conflictos de interés que revelar.    Los
roles de autor se clasificaron utilizando la taxonomía de roles de
colaborador (CRediT; https://credit.niso.org/) de la siguiente
manera:  Edison Achalma:   conceptualización, metodología, análisis
formal, investigación, recursos, curación de
datos, redacción, visualización, supervisión, administración del
proyecto}
\par{La correspondencia relativa a este artículo debe dirigirse a Edison
Achalma, Escuela Profesional de Economía, Universidad Nacional de San
Cristóbal de Huamanga, Portal Independencia N°
57, Ayacucho, AYA 5001, Perú, Email: \href{mailto:elmer.achalma.09@unsch.edu.pe}{elmer.achalma.09@unsch.edu.pe}}
}

\makeatletter
\let\endoldlt\endlongtable
\def\endlongtable{
\hline
\endoldlt
}
\makeatother

\urlstyle{same}



\hypersetup{
  pdftitle={Nombres de variables},
  pdfauthor={Autor}
}
\makeatletter
\@ifpackageloaded{caption}{}{\usepackage{caption}}
\AtBeginDocument{%
\ifdefined\contentsname
  \renewcommand*\contentsname{Tabla de contenidos}
\else
  \newcommand\contentsname{Tabla de contenidos}
\fi
\ifdefined\listfigurename
  \renewcommand*\listfigurename{Lista de Figuras}
\else
  \newcommand\listfigurename{Lista de Figuras}
\fi
\ifdefined\listtablename
  \renewcommand*\listtablename{Lista de Tablas}
\else
  \newcommand\listtablename{Lista de Tablas}
\fi
\ifdefined\figurename
  \renewcommand*\figurename{Figura}
\else
  \newcommand\figurename{Figura}
\fi
\ifdefined\tablename
  \renewcommand*\tablename{Tabla}
\else
  \newcommand\tablename{Tabla}
\fi
}
\@ifpackageloaded{float}{}{\usepackage{float}}
\floatstyle{ruled}
\@ifundefined{c@chapter}{\newfloat{codelisting}{h}{lop}}{\newfloat{codelisting}{h}{lop}[chapter]}
\floatname{codelisting}{Listado}
\newcommand*\listoflistings{\listof{codelisting}{Lista de Listados}}
\makeatother
\makeatletter
\makeatother
\makeatletter
\@ifpackageloaded{caption}{}{\usepackage{caption}}
\@ifpackageloaded{subcaption}{}{\usepackage{subcaption}}
\makeatother
\makeatletter
\@ifpackageloaded{fontawesome5}{}{\usepackage{fontawesome5}}
\makeatother

% From https://tex.stackexchange.com/a/645996/211326
%%% apa7 doesn't want to add appendix section titles in the toc
%%% let's make it do it
\makeatletter
\xpatchcmd{\appendix}
  {\par}
  {\addcontentsline{toc}{section}{\@currentlabelname}\par}
  {}{}
\makeatother

%% Disable longtable counter
%% https://tex.stackexchange.com/a/248395/211326

\usepackage{etoolbox}

\makeatletter
\patchcmd{\LT@caption}
  {\bgroup}
  {\bgroup\global\LTpatch@captiontrue}
  {}{}
\patchcmd{\longtable}
  {\par}
  {\par\global\LTpatch@captionfalse}
  {}{}
\apptocmd{\endlongtable}
  {\ifLTpatch@caption\else\addtocounter{table}{-1}\fi}
  {}{}
\newif\ifLTpatch@caption
\makeatother

\begin{document}

\maketitle


\hypertarget{toc}{}
\tableofcontents
\newpage
\section[Introduction]{Nombres de variables}

\setcounter{secnumdepth}{3}

\setlength\LTleft{0pt}




\subsection{🎯 Motivación}\label{motivaciuxf3n}

Los nombres de variables son la forma en que te comunicas con otros
programadores y contigo mismo. Un buen nombre hace que tu código sea
\textbf{autoexplicativo}, reduciendo la necesidad de comentarios y
facilitando el mantenimiento. Elegir nombres claros no es solo estilo:
es una habilidad fundamental que distingue a un programador novato de
uno profesional.

\subsection{💡 Intuición}\label{intuiciuxf3n}

Piensa en los nombres de variables como las etiquetas de los cajones en
una oficina. Si cada cajón tiene una etiqueta clara (``Facturas'',
``Contratos'', ``Informes''), cualquiera puede encontrar lo que busca.
Si usas etiquetas como ``Cosa1'', ``Cosa2'' o ``Eso'', nadie (ni tú)
sabrá qué hay dentro. En programación pasa igual: \texttt{total\_ventas}
es mucho mejor que \texttt{x} o \texttt{temp}.

\subsection{🧠 Modelo mental}\label{modelo-mental}

Las reglas de nombres en Python son \textbf{estrictas} (el intérprete te
dará error si no las cumples) y \textbf{convencionales} (buenas
prácticas que no son obligatorias pero que mejoran la legibilidad).

\subsubsection{Reglas obligatorias}\label{reglas-obligatorias}

\begin{itemize}
\tightlist
\item
  Solo letras (A-Z, a-z), números (0-9) y guión bajo (\texttt{\_})
\item
  No puede empezar con número
\item
  No puede ser una palabra reservada de Python
\item
  Sensible a mayúsculas y minúsculas
\end{itemize}

\subsubsection{Convenciones}\label{convenciones}

\begin{itemize}
\tightlist
\item
  \texttt{snake\_case} (recomendado por PEP 8)
\item
  Nombres descriptivos
\item
  Constantes en \texttt{MAYUSCULAS}
\item
  Clases en \texttt{PascalCase}
\end{itemize}

\subsection{📐 Desarrollo}\label{desarrollo}

\subsubsection{1. Reglas estrictas (lo que es legal o
ilegal)}\label{reglas-estrictas-lo-que-es-legal-o-ilegal}

\textbf{Nombres legales} ✅:

\begin{Shaded}
\begin{Highlighting}[]
\NormalTok{nombre }\OperatorTok{=} \StringTok{"Ana"}        \CommentTok{\# Letras}
\NormalTok{\_nombre }\OperatorTok{=} \StringTok{"Ana"}       \CommentTok{\# Empieza con guión bajo}
\NormalTok{nombre2 }\OperatorTok{=} \StringTok{"Ana"}       \CommentTok{\# Termina con número}
\NormalTok{nombre\_completo }\OperatorTok{=} \StringTok{"Ana"}  \CommentTok{\# Guión bajo}
\NormalTok{Nombre }\OperatorTok{=} \StringTok{"Ana"}        \CommentTok{\# Mayúscula inicial (pero no recomendado)}
\NormalTok{NOMBRE }\OperatorTok{=} \StringTok{"Ana"}        \CommentTok{\# Todo mayúsculas (para constantes)}
\end{Highlighting}
\end{Shaded}

\textbf{Nombres ilegales} ❌:

\begin{Shaded}
\begin{Highlighting}[]
\DecValTok{2}\ErrorTok{nombre} \OperatorTok{=} \StringTok{"Ana"}       \CommentTok{\# Empieza con número}
\NormalTok{mi}\OperatorTok{{-}}\NormalTok{nombre }\OperatorTok{=} \StringTok{"Ana"}     \CommentTok{\# Contiene guión ({-})}
\NormalTok{mi nombre }\OperatorTok{=} \StringTok{"Ana"}     \CommentTok{\# Contiene espacio}
\ControlFlowTok{if} \OperatorTok{=} \StringTok{"Ana"}            \CommentTok{\# Palabra reservada}
\end{Highlighting}
\end{Shaded}

\subsubsection{2. Convenciones de estilo para nombres
compuestos}\label{convenciones-de-estilo-para-nombres-compuestos}

Cuando un nombre tiene varias palabras, usa una de estas convenciones:

{\def\LTcaptype{none} % do not increment counter
\begin{longtable}[]{@{}
  >{\raggedright\arraybackslash}p{(\linewidth - 4\tabcolsep) * \real{0.1974}}
  >{\raggedright\arraybackslash}p{(\linewidth - 4\tabcolsep) * \real{0.2632}}
  >{\raggedright\arraybackslash}p{(\linewidth - 4\tabcolsep) * \real{0.5395}}@{}}
\toprule\noalign{}
\begin{minipage}[b]{\linewidth}\raggedright
Estilo
\end{minipage} & \begin{minipage}[b]{\linewidth}\raggedright
Ejemplo
\end{minipage} & \begin{minipage}[b]{\linewidth}\raggedright
Uso recomendado
\end{minipage} \\
\midrule\noalign{}
\endhead
\bottomrule\noalign{}
\endlastfoot
\textbf{Snake Case} & \texttt{mi\_nombre\_variable} & Variables y
funciones (PEP 8) \\
\textbf{Camel Case} & \texttt{miNombreVariable} & A veces usado en
JavaScript, no en Python \\
\textbf{Pascal Case} & \texttt{MiNombreVariable} & Clases en Python \\
\end{longtable}
}

\textbf{Recomendación:} usa \texttt{snake\_case} para variables en
Python:

\begin{Shaded}
\begin{Highlighting}[]
\CommentTok{\# ✅ Snake Case (recomendado)}
\NormalTok{precio\_total }\OperatorTok{=} \DecValTok{100}
\NormalTok{nombre\_usuario }\OperatorTok{=} \StringTok{"Carlos"}
\NormalTok{lista\_productos }\OperatorTok{=}\NormalTok{ [}\StringTok{"manzana"}\NormalTok{, }\StringTok{"pera"}\NormalTok{]}

\CommentTok{\# ❌ Camel Case (no recomendado en Python)}
\NormalTok{precioTotal }\OperatorTok{=} \DecValTok{100}
\NormalTok{nombreUsuario }\OperatorTok{=} \StringTok{"Carlos"}
\end{Highlighting}
\end{Shaded}

\subsubsection{3. Palabras reservadas (no
usar)}\label{palabras-reservadas-no-usar}

Estas palabras tienen un significado especial en Python y no pueden
usarse como nombres:

\begin{Shaded}
\begin{Highlighting}[]
\CommentTok{\# No usar como variables}
\CommentTok{\# False, class, finally, is, return, None, continue, for, lambda, try,}
\CommentTok{\# True, def, from, nonlocal, while, and, del, global, not, with,}
\CommentTok{\# as, elif, if, or, yield, assert, else, import, pass, break, except, in, raise}
\end{Highlighting}
\end{Shaded}

\subsubsection{4. Convenciones especiales (por
contexto)}\label{convenciones-especiales-por-contexto}

{\def\LTcaptype{none} % do not increment counter
\begin{longtable}[]{@{}
  >{\raggedright\arraybackslash}p{(\linewidth - 4\tabcolsep) * \real{0.3750}}
  >{\raggedright\arraybackslash}p{(\linewidth - 4\tabcolsep) * \real{0.2361}}
  >{\raggedright\arraybackslash}p{(\linewidth - 4\tabcolsep) * \real{0.3889}}@{}}
\toprule\noalign{}
\begin{minipage}[b]{\linewidth}\raggedright
Convención
\end{minipage} & \begin{minipage}[b]{\linewidth}\raggedright
Ejemplo
\end{minipage} & \begin{minipage}[b]{\linewidth}\raggedright
Significado
\end{minipage} \\
\midrule\noalign{}
\endhead
\bottomrule\noalign{}
\endlastfoot
\textbf{Mayúsculas} & \texttt{PI\ =\ 3.14159} & Constante (no debe
cambiar) \\
\textbf{Guión bajo al inicio} & \texttt{\_interno} & ``Privado'' (por
convención) \\
\textbf{Doble guión bajo} & \texttt{\_\_magico\_\_} & Métodos especiales
de Python \\
\textbf{Minúsculas y guión bajo} & \texttt{variable\_normal} & Variables
estándar \\
\end{longtable}
}

\subsection{🔬 Ejemplos}\label{ejemplos}

\textbf{Mínimo ---} Nombres legales e ilegales:

\begin{Shaded}
\begin{Highlighting}[]
\CommentTok{\# ✅ Legales}
\NormalTok{nombre }\OperatorTok{=} \StringTok{"Ana"}
\NormalTok{\_edad }\OperatorTok{=} \DecValTok{25}
\NormalTok{precio\_producto }\OperatorTok{=} \FloatTok{99.99}
\NormalTok{total\_1 }\OperatorTok{=} \DecValTok{100}

\CommentTok{\# ❌ Ilegales (descomentar da error)}
\CommentTok{\# 1er\_lugar = "primero"     \# Empieza con número}
\CommentTok{\# mi{-}variable = 10          \# Tiene guión}
\CommentTok{\# class = "programacion"    \# Palabra reservada}
\CommentTok{\# mi variable = "texto"     \# Tiene espacio}
\end{Highlighting}
\end{Shaded}

\textbf{Progresivo ---} Nombres descriptivos vs.~nombres confusos:

\begin{Shaded}
\begin{Highlighting}[]
\CommentTok{\# ❌ Malos nombres (difíciles de entender)}
\NormalTok{a }\OperatorTok{=} \FloatTok{12.5}
\NormalTok{b }\OperatorTok{=} \StringTok{"Juan"}
\NormalTok{c }\OperatorTok{=}\NormalTok{ a }\OperatorTok{*} \FloatTok{1.21}
\BuiltInTok{print}\NormalTok{(c)}

\CommentTok{\# ✅ Buenos nombres (autoexplicativos)}
\NormalTok{precio\_base }\OperatorTok{=} \FloatTok{12.5}
\NormalTok{nombre\_cliente }\OperatorTok{=} \StringTok{"Juan"}
\NormalTok{precio\_con\_iva }\OperatorTok{=}\NormalTok{ precio\_base }\OperatorTok{*} \FloatTok{1.21}
\BuiltInTok{print}\NormalTok{(precio\_con\_iva)}
\end{Highlighting}
\end{Shaded}

\textbf{Real ---} Sistema de gestión de inventario:

\begin{Shaded}
\begin{Highlighting}[]
\CommentTok{\# Configuración del sistema}
\NormalTok{LIMITE\_INVENTARIO }\OperatorTok{=} \DecValTok{100}          \CommentTok{\# Constante (mayúsculas)}
\NormalTok{ruta\_archivo }\OperatorTok{=} \StringTok{"/data/inventario.csv"}  \CommentTok{\# Variable normal}
\NormalTok{\_contador\_interno }\OperatorTok{=} \DecValTok{0}            \CommentTok{\# Variable "privada"}

\CommentTok{\# Datos de productos}
\NormalTok{producto\_nombre }\OperatorTok{=} \StringTok{"Laptop"}
\NormalTok{producto\_precio }\OperatorTok{=} \FloatTok{850.50}
\NormalTok{producto\_stock }\OperatorTok{=} \DecValTok{15}
\NormalTok{producto\_activo }\OperatorTok{=} \VariableTok{True}

\CommentTok{\# Mensaje de estado}
\BuiltInTok{print}\NormalTok{(}\StringTok{"Producto:"}\NormalTok{, producto\_nombre)}
\BuiltInTok{print}\NormalTok{(}\StringTok{"Stock:"}\NormalTok{, producto\_stock, }\StringTok{"unidades"}\NormalTok{)}
\BuiltInTok{print}\NormalTok{(}\StringTok{"Stock bajo:"}\NormalTok{, producto\_stock }\OperatorTok{\textless{}} \DecValTok{10}\NormalTok{)  }\CommentTok{\# Resultado: False}
\end{Highlighting}
\end{Shaded}

\subsection{⚠️ Errores comunes}\label{errores-comunes}

{\def\LTcaptype{none} % do not increment counter
\begin{longtable}[]{@{}
  >{\raggedright\arraybackslash}p{(\linewidth - 4\tabcolsep) * \real{0.2581}}
  >{\raggedright\arraybackslash}p{(\linewidth - 4\tabcolsep) * \real{0.2903}}
  >{\raggedright\arraybackslash}p{(\linewidth - 4\tabcolsep) * \real{0.4516}}@{}}
\toprule\noalign{}
\begin{minipage}[b]{\linewidth}\raggedright
Error
\end{minipage} & \begin{minipage}[b]{\linewidth}\raggedright
Síntoma
\end{minipage} & \begin{minipage}[b]{\linewidth}\raggedright
Cómo evitarlo
\end{minipage} \\
\midrule\noalign{}
\endhead
\bottomrule\noalign{}
\endlastfoot
\textbf{Espacio en nombre} & \texttt{SyntaxError:\ invalid\ syntax} &
Usa \texttt{\_} en lugar de espacio \\
\textbf{Número al inicio} & \texttt{SyntaxError:\ invalid\ syntax} &
Empieza con letra o \texttt{\_} \\
\textbf{Palabra reservada} & \texttt{SyntaxError:\ invalid\ syntax} &
Consulta la lista de palabras reservadas \\
\textbf{Guión en lugar de guión bajo} &
\texttt{SyntaxError:\ invalid\ syntax} & Usa \texttt{\_} no
\texttt{-} \\
\textbf{Mayúsculas incorrectas} &
\texttt{NameError:\ name\ \textquotesingle{}x\textquotesingle{}\ is\ not\ defined}
& Python es sensible a mayúsculas: \texttt{Variable} ≠
\texttt{variable} \\
\end{longtable}
}

\subsection{✅ Buenas prácticas}\label{buenas-pruxe1cticas}

\begin{enumerate}
\def\labelenumi{\arabic{enumi}.}
\item
  \textbf{Nombres descriptivos pero concisos}:

\begin{Shaded}
\begin{Highlighting}[]
\CommentTok{\# ❌ Demasiado corto}
\NormalTok{a }\OperatorTok{=} \DecValTok{100}

\CommentTok{\# ❌ Demasiado largo}
\NormalTok{precio\_total\_de\_la\_compra\_sin\_iva }\OperatorTok{=} \DecValTok{100}

\CommentTok{\# ✅ Justo}
\NormalTok{precio\_base }\OperatorTok{=} \DecValTok{100}
\end{Highlighting}
\end{Shaded}
\item
  \textbf{Usa verbos para funciones, sustantivos para variables}:

\begin{Shaded}
\begin{Highlighting}[]
\CommentTok{\# Función: acción (verbo)}
\KeywordTok{def}\NormalTok{ calcular\_total():}
    \ControlFlowTok{pass}

\CommentTok{\# Variable: dato (sustantivo)}
\NormalTok{total }\OperatorTok{=} \DecValTok{100}
\end{Highlighting}
\end{Shaded}
\item
  \textbf{Sé consistente}:

\begin{Shaded}
\begin{Highlighting}[]
\CommentTok{\# ❌ Inconsistente}
\NormalTok{nombre\_usuario }\OperatorTok{=} \StringTok{"Ana"}
\NormalTok{user\_email }\OperatorTok{=} \StringTok{"ana@mail.com"}

\CommentTok{\# ✅ Consistente}
\NormalTok{nombre\_usuario }\OperatorTok{=} \StringTok{"Ana"}
\NormalTok{email\_usuario }\OperatorTok{=} \StringTok{"ana@mail.com"}
\end{Highlighting}
\end{Shaded}
\item
  \textbf{Evita abreviaturas no estándar}:

\begin{Shaded}
\begin{Highlighting}[]
\CommentTok{\# ❌ Confuso}
\NormalTok{prv }\OperatorTok{=} \StringTok{"Programación"}

\CommentTok{\# ✅ Claro}
\NormalTok{curso }\OperatorTok{=} \StringTok{"Programación"}
\end{Highlighting}
\end{Shaded}
\item
  \textbf{Usa nombres que tengan sentido en el contexto del programa}.
\end{enumerate}

\subsection{📝 Ejercicio}\label{ejercicio}

\textbf{Básico} --- Corrige nombres ilegales

Los siguientes nombres de variables tienen errores. Corrígelos para que
sean válidos y sigan las convenciones de Python:

\begin{Shaded}
\begin{Highlighting}[]
\CommentTok{\# Corrige estos nombres ilegales/malos}
\DecValTok{1}\ErrorTok{er\_trimestre} \OperatorTok{=} \DecValTok{2024}
\NormalTok{total}\OperatorTok{{-}}\NormalTok{ventas }\OperatorTok{=} \DecValTok{15000}
\NormalTok{mi variable }\OperatorTok{=} \StringTok{"casa"}
\KeywordTok{class} \OperatorTok{=} \StringTok{"Python"}
\NormalTok{NombreDeUsuario }\OperatorTok{=} \StringTok{"ana"}
\NormalTok{precioTotal }\OperatorTok{=} \FloatTok{99.99}
\end{Highlighting}
\end{Shaded}

\begin{center}\rule{0.5\linewidth}{0.5pt}\end{center}

\textbf{Solución (no mires hasta haberlo intentado):}

\begin{Shaded}
\begin{Highlighting}[]
\CommentTok{\# Versión corregida (snake case)}
\NormalTok{primer\_trimestre }\OperatorTok{=} \DecValTok{2024}
\NormalTok{total\_ventas }\OperatorTok{=} \DecValTok{15000}
\NormalTok{mi\_variable }\OperatorTok{=} \StringTok{"casa"}
\NormalTok{clase }\OperatorTok{=} \StringTok{"Python"}  \CommentTok{\# O mejor: nombre\_clase}
\NormalTok{nombre\_usuario }\OperatorTok{=} \StringTok{"ana"}
\NormalTok{precio\_total }\OperatorTok{=} \FloatTok{99.99}
\end{Highlighting}
\end{Shaded}

\subsection{🗂️ Síntesis}\label{suxedntesis}

\begin{itemize}
\tightlist
\item
  \textbf{Reglas obligatorias}: letras, números, \texttt{\_}; no empezar
  con número; sin espacios ni \texttt{-}; sin palabras reservadas.
\item
  \textbf{Sensible a mayúsculas}: \texttt{nombre} ≠ \texttt{Nombre} ≠
  \texttt{NOMBRE}.
\item
  \textbf{Snake Case}: convención recomendada para variables
  (\texttt{mi\_variable}).
\item
  \textbf{Descriptivos}: usa nombres que revelen el propósito de la
  variable.
\item
  \textbf{Constantes en mayúsculas}: \texttt{LIMITE\ =\ 100}.
\item
  \textbf{Palabras reservadas}: no usar \texttt{if}, \texttt{for},
  \texttt{class}, \texttt{while}, etc.
\end{itemize}

\subsection{🔗 Qué sigue}\label{quuxe9-sigue}

\begin{itemize}
\tightlist
\item
  \textbf{Relacionado:} variables (asignación, tipos), palabras
  reservadas
\item
  \textbf{Siguiente tema:} Tipos de datos --- aprenderás a manejar
  diferentes tipos de datos y sus propiedades.
\end{itemize}

\section{Publicaciones Similares}\label{publicaciones-similares}

Si te interesó este artículo, te recomendamos que explores otros blogs y
recursos relacionados que pueden ampliar tus conocimientos. Aquí te dejo
algunas sugerencias:

\begin{enumerate}
\def\labelenumi{\arabic{enumi}.}
\tightlist
\item
  \href{https://numerus-scriptum.netlify.app/python/2020-06-19-instalacion-de-anaconda/index.pdf}{\faIcon{file-pdf}}
  \href{https://numerus-scriptum.netlify.app/python/2020-06-19-instalacion-de-anaconda}{Instalacion
  De Anaconda}
\item
  \href{https://numerus-scriptum.netlify.app/python/2020-06-20-configurar-entorno-virtual-python-anaconda/index.pdf}{\faIcon{file-pdf}}
  \href{https://numerus-scriptum.netlify.app/python/2020-06-20-configurar-entorno-virtual-python-anaconda}{Configurar
  Entorno Virtual Python Anaconda}
\item
  \href{https://numerus-scriptum.netlify.app/python/2021-04-17-01-introducion-a-la-programacion-con-python/index.pdf}{\faIcon{file-pdf}}
  \href{https://numerus-scriptum.netlify.app/python/2021-04-17-01-introducion-a-la-programacion-con-python}{01
  Introducion A La Programacion Con Python}
\item
  \href{https://numerus-scriptum.netlify.app/python/2021-05-31-02-variables-expresiones-y-statements-con-python/index.pdf}{\faIcon{file-pdf}}
  \href{https://numerus-scriptum.netlify.app/python/2021-05-31-02-variables-expresiones-y-statements-con-python}{02
  Variables Expresiones Y Statements Con Python}
\item
  \href{https://numerus-scriptum.netlify.app/python/2021-06-07-03-objetos-de-python/index.pdf}{\faIcon{file-pdf}}
  \href{https://numerus-scriptum.netlify.app/python/2021-06-07-03-objetos-de-python}{03
  Objetos De Python}
\item
  \href{https://numerus-scriptum.netlify.app/python/2021-06-14-04-ejecucion-condicional-con-python/index.pdf}{\faIcon{file-pdf}}
  \href{https://numerus-scriptum.netlify.app/python/2021-06-14-04-ejecucion-condicional-con-python}{04
  Ejecucion Condicional Con Python}
\item
  \href{https://numerus-scriptum.netlify.app/python/2021-06-21-05-iteraciones-con-python/index.pdf}{\faIcon{file-pdf}}
  \href{https://numerus-scriptum.netlify.app/python/2021-06-21-05-iteraciones-con-python}{05
  Iteraciones Con Python}
\item
  \href{https://numerus-scriptum.netlify.app/python/2021-08-16-06-funciones-con-python/index.pdf}{\faIcon{file-pdf}}
  \href{https://numerus-scriptum.netlify.app/python/2021-08-16-06-funciones-con-python}{06
  Funciones Con Python}
\item
  \href{https://numerus-scriptum.netlify.app/python/2021-08-23-07-dataframes-con-python/index.pdf}{\faIcon{file-pdf}}
  \href{https://numerus-scriptum.netlify.app/python/2021-08-23-07-dataframes-con-python}{07
  Dataframes Con Python}
\item
  \href{https://numerus-scriptum.netlify.app/python/2021-11-29-08-prediccion-y-metrica-de-performance-con-python/index.pdf}{\faIcon{file-pdf}}
  \href{https://numerus-scriptum.netlify.app/python/2021-11-29-08-prediccion-y-metrica-de-performance-con-python}{08
  Prediccion Y Metrica De Performance Con Python}
\item
  \href{https://numerus-scriptum.netlify.app/python/2021-12-06-09-metodos-de-machine-learning-para-clasificacion-con-python/index.pdf}{\faIcon{file-pdf}}
  \href{https://numerus-scriptum.netlify.app/python/2021-12-06-09-metodos-de-machine-learning-para-clasificacion-con-python}{09
  Metodos De Machine Learning Para Clasificacion Con Python}
\item
  \href{https://numerus-scriptum.netlify.app/python/2021-12-13-10-metodos-de-machine-learning-para-regresion-con-python/index.pdf}{\faIcon{file-pdf}}
  \href{https://numerus-scriptum.netlify.app/python/2021-12-13-10-metodos-de-machine-learning-para-regresion-con-python}{10
  Metodos De Machine Learning Para Regresion Con Python}
\item
  \href{https://numerus-scriptum.netlify.app/python/2025-05-10-visualizacion-de-datos-con-python/index.pdf}{\faIcon{file-pdf}}
  \href{https://numerus-scriptum.netlify.app/python/2025-05-10-visualizacion-de-datos-con-python}{Visualizacion
  De Datos Con Python}
\end{enumerate}

Esperamos que encuentres estas publicaciones igualmente interesantes y
útiles. ¡Disfruta de la lectura!






\end{document}
